% TOP_PROBLEM: {"id": 12, "title": "Goldbach's Conjecture", "category_id": 1, "category": {"id": 1, "name": "number_theory", "display_name": "Number Theory", "description": "Properties of integers, prime numbers, Diophantine equations.", "slug": "number-theory", "order_index": 1, "created_at": "2026-07-31T15:26:25.670Z"}, "status": "open"}
% FIRST_POSTED: 2026-09-27
% !TeX program = lualatex
% ATTEMPT_STATUS: unresolved
% Research continuation by GPT_6_Astra_Ultra, 27 September 2026.
% Catalogue statements and existing sources preserved verbatim.
\documentclass[11pt]{article}
\usepackage[margin=25mm]{geometry}
\usepackage{fontspec}
\setmainfont{Latin Modern Roman}
\usepackage{amsmath,amssymb,mathtools}
\usepackage{unicode-math}
\setmathfont{Latin Modern Math}
\usepackage{xurl,hyperref}
\hypersetup{colorlinks=true,urlcolor=blue}
\setcounter{secnumdepth}{0}
\setlength{\parindent}{0pt}
\setlength{\parskip}{5pt}
\setlength{\emergencystretch}{3em}
\begin{document}
\section*{\texorpdfstring{Goldbach's Conjecture}{Goldbach's Conjecture}}\label{12}
\subsection{Definitions and mathematical statement}
With $\mathcal P$ denoting the positive primes, the statement is
\[
 \forall n\in2{\mathbb{Z}},\quad n>2\quad\Longrightarrow\quad
 \exists p,q\in\mathcal P\text{ with }n=p+q.
\]
The primes are not required to be distinct, so $4=2+2$ is allowed. Every even integer is quantified, not merely all sufficiently large integers or a density-one subset.
\par\smallskip\textit{Formulation and concept references:} \href{https://t5k.org/notes/conjectures/}{[S1]}.

\subsection{Short English statement}
Can every even integer larger than two be written as the sum of two primes?
\subsection{Research attempt}

\paragraph{Scope and source check.}
The target includes every even integer greater than two, allows equal
primes, and therefore includes $4=2+2$. An asymptotic argument for large
integers needs an explicit finite-range completion before it proves this
formulation. Oliveira e Silva--Herzog--Pardi [E1] report verification
through $4\cdot10^{18}$; that published computation is not rerun here.
It is a finite theorem, not a replacement for control of the unbounded tail.

The recent catalogue source [S2], Bhowmik--Grimmelt's August 2026 revision,
is largely a survey of exceptional-set methods, with additional explicit
formula work. It does not assert that the exceptional set is empty.
Li--Liu [S3], version 2, state a representation
$N=p+rq$, with $p,q$ prime, $r=1$ or prime, and $r\le q^{0.9}$ for
sufficiently large even $N$. This still permits $r>1$, so its stated
conclusion is weaker than binary Goldbach. The manuscript statements
were inspected; their full proofs were not independently audited.
Neither claim is used as a premise of the elementary arguments below.

The selected route asks exactly what the circle method would have to
prove for one prescribed even integer, and compares that with what an
average-square error bound actually guarantees. The goal is to identify
the first missing uniform estimate, not to infer pointwise positivity
from an exceptional set of density zero.

\paragraph{An exact Fourier representation without prime powers.}
For an integer $X\ge4$, put $e(t)=\exp(2\pi i t)$ and define
\[
 S_X(\alpha)=\sum_{p\le X}(\log p)e(p\alpha),\qquad
 R(N)=\sum_{\substack{p+q=N\\p,q\text{ prime}}}\log p\log q.
\]
The sum counts ordered pairs. For $4\le N\le X$, orthogonality gives
\begin{equation}\label{ra:25:fourier}
 R(N)=\int_0^1 S_X(\alpha)^2e(-N\alpha)\,d\alpha.
\end{equation}
This follows by expanding a finite double sum and using
$\int_0^1e(m\alpha)\,d\alpha=1$ for $m=0$, zero otherwise.
All weights are positive, so $R(N)>0$ is exactly the desired existence
statement. Proper prime powers are absent by definition, and no
subtraction estimate for them is hidden in this equivalence.

The square in \eqref{ra:25:fourier} is not an absolute square.
Replacing it by $|S_X|^2$ counts prime \emph{differences} $p-q=N$.
For example with $X=N=4$, the difference count is zero, whereas
$R(4)=(\log2)^2>0$. Positivity of $|S_X|^2$ therefore cannot be used
to prove positivity of the required Fourier coefficient. The actual
coefficient is real, but substantial complex cancellation occurs inside
the integral.

\paragraph{Local factors give a positive predicted main term.}
For an odd prime $r$, count residue pairs $(a,b)$ with $a+b=N$ modulo
$r$ and both coordinates nonzero. The count is $r-1$ when $r\mid N$
and $r-2$ otherwise: choose $a$, excluding zero and, in the second case,
the distinct residue $N$. Relative to independent unit densities, the
local multiplier is
\[
 \frac{r\,\#\{a\bmod r:a\ne0,\ N-a\ne0\}}{(r-1)^2}.
\]
For even $N$, the factor at two is two. Consequently the formal
singular series can be written
\begin{equation}\label{ra:25:series}
 \mathfrak S_G(N)=2C_2
      \prod_{\substack{r\mid N\\r>2}}\frac{r-1}{r-2},
 \qquad C_2=\prod_{r>2}\left(1-\frac1{(r-1)^2}\right).
\end{equation}
The infinite product is strictly positive: the factors lie in $(0,1)$,
and the sum of their deficits converges by comparison with
$\sum_{n\ge2}n^{-2}$. For instance $-\log(1-u)\le2u$ for
$0\le u\le1/2$ bounds the logarithmic tail. Thus
$\mathfrak S_G(N)\ge c_0:=2C_2>0$ for every even $N$.

This derivation is an exact local calculation, not a proof that
$R(N)\sim\mathfrak S_G(N)N$. CRT similarly multiplies the residue counts
over a fixed squarefree modulus, but actual primality is a global
restriction. The positive lower bound for the proposed main term is
useful only if an accompanying analytic error can be controlled.

The small prime two requires a separate boundary check. If $N>4$ is
even and one summand were two, the other would be the even integer
$N-2>2$ and could not be prime. Thus every representation for even
$N>4$ uses two odd primes, consistent with the local unit calculation
at two. The exceptional representation of four is handled by the exact
Fourier identity and the finite check, not by an asymptotic local model.
Similarly, ordered counting never doubles a diagonal pair: when $N/2$
is prime it contributes one term of weight $(\log(N/2))^2$.
These conventions ensure that the positive coefficient being targeted
matches the precise existence statement at both the boundary and large
values of $N$.

\paragraph{A precise pointwise sufficient estimate.}
Choose $Q=(\log X)^B$ with fixed $B>0$ and define major arcs as the
union, modulo one, of
\[
 |\alpha-a/q|\le Q/(qX),\qquad
 1\le q\le Q,\quad 0\le a<q,\quad(a,q)=1.
\]
Let this union be $\mathfrak M_X$, with complement $\mathfrak m_X$,
and set
\[
 A_X(N)=\int_{\mathfrak M_X}S_X(\alpha)^2e(-N\alpha)\,d\alpha,
 \quad I_X(N)=\int_{\mathfrak m_X}S_X(\alpha)^2e(-N\alpha)\,d\alpha.
\]
The sets are invariant under $\alpha\mapsto-\alpha$, making both
integrals real. They sum to $R(N)$ exactly. Their definition is valid
even when arcs overlap; for large $X$ the usual disjointness also follows
from $X>2Q^3$ and separation of reduced fractions.

Suppose, for some explicit $X_0$ and all $X\ge X_0$, one could establish
uniformly for even $X/2\le N\le X$ that
\begin{equation}\label{ra:25:targets}
 |A_X(N)-\mathfrak S_G(N)N|\le c_0X/8,
 \qquad |I_X(N)|<c_0X/8.
\end{equation}
Then $R(N)>c_0X/4>0$. Dyadic intervals would cover every sufficiently
large even integer. An explicit check of the remaining interval
$4\le N<X_0$ would finish the exact catalogue statement. This implication
is proved by the displayed estimates; the estimates themselves are not
proved here. Classical major-arc analysis and the role of exceptional
sets are discussed in [S2]. In this notebook no theorem with different
uniformity or an implicit exceptional set is substituted for
\eqref{ra:25:targets}.

\paragraph{What a mean-square estimate really provides.}
The elementary Bessel inequality already quantifies the loss when one
passes from a single integer to most integers. Fix $X$ and suppose a
major-arc lower bound $A_X(N)\ge aX$ is available on its even half-interval,
for a fixed $a>0$. If such an $N$ has no representation, then
$I_X(N)=-A_X(N)$, so $|I_X(N)|\ge aX$.
Let $U=\sup_{\alpha\in\mathfrak m_X}|S_X(\alpha)|$.
Orthogonality and Bessel give
\begin{align*}
 \sum_{N\in\mathbb Z}|I_X(N)|^2
 &=\int_{\mathfrak m_X}|S_X(\alpha)|^4\,d\alpha\\
 &\le U^2\int_0^1|S_X(\alpha)|^2\,d\alpha
   =U^2\sum_{p\le X}(\log p)^2
   \le U^2 X\log^2 X.
\end{align*}
The first equality is Parseval for the square-integrable function
$1_{\mathfrak m_X}S_X^2$; restricting the sum to even $N$ would give
an inequality, which is all we need. Therefore the number $E_X$ of
exceptional even integers in that interval satisfies
\begin{equation}\label{ra:25:exceptional}
 E_X\le\frac{U^2\log^2X}{a^2X}.
\end{equation}
This derivation uses only the elementary bound on the number of primes;
sharper familiar second-moment estimates are unnecessary for the point.

For example, even a bound $U\le X/(\log X)^A$ with fixed $A>1$ would
give only $E_X\ll_a X/(\log X)^{2A-2}$. Its density tends to zero,
but the bound grows and does not force the integer $E_X$ to vanish.
Applying Cauchy--Schwarz directly to one coefficient instead gives
$|I_X(N)|\le U\sqrt X\log X$, which is useful for positivity only
with a far stronger supremum bound. Neither estimate captures the
specific oscillatory cancellation of $e(-N\alpha)$ needed pointwise.

\paragraph{A sharp stress test of the average-to-pointwise step.}
Consider an auxiliary array with predicted main value $aX$ at every
even index in $[X/2,X]$. Let its actual value equal that prediction
except at one index, where it is zero. All values remain nonnegative,
and the sum of squared errors is exactly $a^2X^2=o(X^3)$.
Thus even a mean-square error small relative to the aggregate natural
scale $X^3$ is compatible with a missing representation. Choosing one
such exceptional index in each dyadic interval makes infinitely many
exceptions of density zero. These arrays are not prime counts and do
not disprove Goldbach; they refute the logical inference from the stated
norm bound to universal positivity.

An average theorem supplemented by finitely many direct checks has the
same issue: the hypothetical sparse exceptional indices may all lie
beyond the checked range. To use integrality, the final upper bound on
$E_X$ must actually be strictly below one for every remaining interval,
or a separate argument must remove each possible exceptional integer.
Calling a positive real upper bound ``small'' cannot replace this step.

\paragraph{Exact finite checks and counting conventions.}
The local script builds the prime indicator by an integer sieve and
computes its ordinary polynomial convolution. If $g(N)$ is the number
of unordered pairs with $p\le q$, the ordered count satisfies
\[
 r(N)=2g(N)-1_{N/2\text{ prime}}.
\]
It checks this identity and verifies every even $4\le N\le10000$ has
a representation. Selected values are
\[
\begin{array}{r|r|r}
N&r(N)&g(N)\\\hline
4&1&1\\10&3&2\\100&12&6\\1000&56&28\\10000&254&127
\end{array}
\]
The diagonal correction is essential at four and ten. The same script
enumerates all residues modulo $210$ and checks the product of the local
counts $r-1$ or $r-2$ in \eqref{ra:25:series}; for example there are
twenty admissible ordered residue pairs for $N=100$ modulo $210$.
These exact local and finite computations do not estimate either arc
integral. Reproduction files are \texttt{checks.py} and
\texttt{results.json} in
\texttt{\_Local/top\_problems/continuation\_15\_25\_2026\_09\_27/analytic/}.

\paragraph{Remaining gap and outcome.}
The work proves the exact prime-weighted Fourier identity, the positive
local-series bound, and the exceptional-set inequality with all
normalizations visible. The first unresolved analytic step is a bound
for the individual minor-arc coefficient strong enough to preserve
positivity, paired with a major-arc estimate of the stated uniformity.
Even proving eventual positivity would leave an effective threshold and
the finite interval below it to address the full target.

\textbf{Outcome: UNRESOLVED.} No pointwise estimate of the required
strength and no counterexample even integer were obtained. The explicit
norm obstruction explains why an almost-all conclusion cannot be promoted
to the requested theorem. The process used direct derivation, primary
source inspection, and exact finite checks; it does not claim that the
open problem has been solved or independently certified.

\subsection{Sources}
\begin{itemize}
\item[S1]C. K. Caldwell, "Prime Conjectures and Open Questions," The PrimePages, https://t5k.org/notes/conjectures/.
\url{https://t5k.org/notes/conjectures/}.
\textit{Formulation source.}
\item[S2]The exceptional set of the Goldbach problem.
\url{https://arxiv.org/html/2607.27282v2}.
\textit{Further reference; not a proof endorsement.}
\item[S3]Theorem (1+1.9) on the Goldbach Conjecture.
\url{https://arxiv.org/abs/2606.05224}.
\textit{Further reference; not a proof endorsement.}
\item[S4]Restricted Goldbach Sums in Arithmetic Progressions: Local Obstructions, an Explicit Almost-All Framework, and a General-Modulus Extension.
\url{https://www.preprints.org/manuscript/202607.2077}.
\textit{Further reference; not a proof endorsement.}

\item[E1]Tom\'as Oliveira e Silva, Siegfried Herzog, and Silvio Pardi,
\emph{Empirical verification of the even Goldbach conjecture and computation
of prime gaps up to $4\cdot10^{18}$}, Mathematics of Computation 83
(2014), no. 288, 2033--2060, DOI: 10.1090/S0025-5718-2013-02787-1.
\url{https://www.ams.org/mcom/2014-83-288/S0025-5718-2013-02787-1/S0025-5718-2013-02787-1.pdf}.
\textit{Primary computational reference; that large computation was not
rerun in this notebook. Consulted 27 September 2026.}
\item[E2]Gautami Bhowmik and Lasse Grimmelt,
\emph{The exceptional set of the Goldbach problem},
arXiv:2607.27282v2 (13 August 2026), especially Sections 1 and 4--7.
\url{https://arxiv.org/html/2607.27282v2}.
\textit{Version-specific identification of [S2]: survey and new explicit-formula
work, not an assertion that every even integer is represented.
Consulted 27 September 2026.}
\item[E3]Jiamin Li and Jianya Liu,
\emph{Theorem $(1+1.9)$ on the Goldbach Conjecture},
arXiv:2606.05224v2 (8 August 2026).
\url{https://arxiv.org/abs/2606.05224v2}.
\textit{Version-specific identification of [S3]: stated representation permits
a semiprime summand. Statement inspected; complete proof not independently
audited or used here. Consulted 27 September 2026.}
\end{itemize}
\end{document}
